% preamble.tex
\usepackage{luatexja}
\usepackage{luatexja-ruby}
\usepackage[match,deluxe]{luatexja-fontspec}

% === Fuentes japonesas (Windows 10 fonts) ===
\setmainjfont{Yu Gothic}[BoldFont=Yu Gothic Bold, ItalicFont=Yu Gothic, BoldItalicFont=Yu Gothic Bold]
\setsansjfont{Yu Gothic}

% === Fuentes latinas ===
\setmainfont{Georgia}
\setsansfont{Calibri}

% === Emoji ===
% Instalar via MiKTeX Package Manager: "emoji"
\usepackage{emoji}
% Si emoji package no disponible, comentar la línea anterior
% y usar: \newcommand{\emoji}[1]{#1}

% === Colores ===
\usepackage{xcolor}
\definecolor{jpred}{HTML}{BC002D}
\definecolor{jpblue}{HTML}{1A3A6B}
\definecolor{jpgold}{HTML}{C8A951}
\definecolor{sakura}{HTML}{FFB7C5}
\definecolor{lightgray}{HTML}{F7F7F7}
\definecolor{darkgray}{HTML}{555555}
\definecolor{softgreen}{HTML}{D4EDDA}

% === Diseño de página ===
\usepackage[a4paper, top=2.5cm, bottom=2.5cm, left=3cm, right=3cm]{geometry}

% === Tablas ===
\usepackage{booktabs}
\usepackage{tabularx}
\usepackage{array}
\usepackage{longtable}

% === tcolorbox ===
\usepackage[most]{tcolorbox}
\tcbuselibrary{skins, breakable}

\tcbset{
  vocabbox/.style={
    enhanced, breakable,
    colback=lightgray, colframe=jpblue,
    coltitle=white, fonttitle=\bfseries\large,
    attach boxed title to top left={yshift=-2mm, xshift=5mm},
    boxed title style={colback=jpblue, rounded corners},
    rounded corners
  },
  grammarbox/.style={
    enhanced, breakable,
    colback=white, colframe=jpred,
    coltitle=white, fonttitle=\bfseries\large,
    attach boxed title to top left={yshift=-2mm, xshift=5mm},
    boxed title style={colback=jpred, rounded corners},
    rounded corners
  },
  phrasebox/.style={
    enhanced,
    colback=sakura!15, colframe=sakura!80!black,
    rounded corners, fonttitle=\bfseries
  },
  kanjibox/.style={
    enhanced, breakable,
    colback=jpgold!8, colframe=jpgold!70!black,
    coltitle=white, fonttitle=\bfseries\large,
    attach boxed title to top left={yshift=-2mm, xshift=5mm},
    boxed title style={colback=jpgold!70!black, rounded corners},
    rounded corners
  },
  ejerciciobox/.style={
    enhanced,
    colback=softgreen!50, colframe=softgreen!80!black,
    rounded corners, fonttitle=\bfseries,
    coltitle=darkgray
  }
}

% === Hyperref ===
\usepackage{hyperref}
\hypersetup{
  colorlinks=true,
  linkcolor=jpblue,
  urlcolor=jpred,
  pdftitle={Apuntes de Japonés N5-N4},
  pdfauthor={Aldo ZM}
}

% === Títulos ===
\usepackage{titlesec}

\titleformat{\chapter}[display]
  {\normalfont\huge\bfseries\color{jpred}}{}
  {0pt}{\huge}
  [\vspace{0.3em}\color{jpred}\hrule height 1.5pt\vspace{0.3em}]

\titleformat{\section}
  {\large\bfseries\color{jpblue}}{}
  {0em}{}
  [\vspace{0.1em}\color{jpblue!40}\hrule height 0.4pt]

\titleformat{\subsection}
  {\normalsize\bfseries\color{darkgray}}{}
  {0em}{}

% =============================================
% COMANDOS CUSTOM
% =============================================

% Encabezado visual de unidad
% #1=número, #2=kanji título, #3=título español, #4=tema comunicativo
\newcommand{\unidadheader}[4]{%
  \begin{tcolorbox}[
    enhanced, colback=jpred!5, colframe=jpred,
    coltitle=white, fonttitle=\bfseries\Large,
    title={🎌 Unidad #1 \enspace \ruby{#2}{} \enspace #3},
    attach boxed title to top center={yshift=-3mm},
    boxed title style={colback=jpred, rounded corners}
  ]
  \textit{\textcolor{jpblue}{🗣️ Tema comunicativo: #4}}
  \end{tcolorbox}
  \vspace{0.5em}
}

% Entrada de vocabulario
% #1=kanji+kana, #2=furigana, #3=español, #4=tipo gramatical
\newcommand{\vocabitem}[4]{%
  \ruby{#1}{#2} & #3 & \textcolor{darkgray}{\small\textit{#4}} \\[4pt]
}

% Patrón gramatical
% #1=nombre patrón, #2=fórmula, #3=explicación+ejemplos
\newcommand{\patron}[3]{%
  \begin{tcolorbox}[grammarbox, title={⚙️ #1}]
  \begin{center}{\Large\textbf{#2}}\end{center}
  \vspace{0.3em}
  #3
  \end{tcolorbox}
}

% Ejemplo dentro de patrón
% #1=japonés (con ruby), #2=español
\newcommand{\ej}[2]{%
  \par\vspace{0.2em}
  \hspace{1em}▶ #1 \hfill {\small\textit{#2}}
  \vspace{0.1em}
}

% Frase clave
% #1=japonés, #2=español
\newcommand{\frase}[2]{%
  \begin{tcolorbox}[phrasebox]
  \textbf{#1}\\[2pt]
  {\small\textit{#2}}
  \end{tcolorbox}
  \vspace{0.2em}
}

% Entrada de kanji en tabla
% #1=kanji, #2=lecturas on, #3=lecturas kun, #4=significado, #5=vocab ejemplo (con ruby)
\newcommand{\kanjirow}[5]{%
  {\LARGE #1} & #2 & #3 & #4 & #5 \\[4pt]
}
